\documentclass[../main.tex]{subfiles}

\begin{document}

\chapter{赋范线性空间的对偶}\label{chapter:ch14-对偶}

对于一个赋范线性空间 $X$, 我们记 $X$ 上的连续线性实值函数的赋范线性空间为 $X^{*}$, 且称它为 $X$ 的对偶空间. 在本章以及下一章, 我们探讨从 $X\times X^{*}$ 到 $\mathbf{R}$ 的映射定义为
\begin{align*}
\text{对所有 }x\in X,\ \psi\in X^{*},\quad(x,\psi)\mapsto\psi(x)
\end{align*}
的映射性质, 以找出 Banach 空间的分析的、几何的、拓扑的性质. 这一探讨的出发点是 Hahn-Banach 定理. 这是关于未赋范的线性空间的子空间上的某种线性泛函到全空间上的线性泛函的延拓定理. 该定理的初等性质给出了这样的灵活性. 本章我们用它导出线性泛函的三个性质: (i)对赋范线性空间 $X$, $X$ 的子空间上的任何有界线性泛函可被延拓为整个 $X$ 上的有界线性泛函, 而它的范数不增加; (ii)对局部凸拓扑向量空间 $X$, $X$ 的任何闭凸子集与该子集外的点可被闭超平面分离; (iii)对自反的 Banach 空间 $X$, $X$ 中的任何有界序列有弱收敛的子序列.

\subfile{Chapter-14/section-01.tex}

\subfile{Chapter-14/section-02.tex}

\subfile{Chapter-14/section-03.tex}

\subfile{Chapter-14/section-04.tex}



\end{document}